\documentclass[border=6pt]{standalone}
\usepackage{fontspec}
\usepackage{xeCJK}
\usepackage{xcolor}
\usepackage{tikz}
\usetikzlibrary{
    arrows.meta,
    backgrounds,
    calc,
    decorations.pathreplacing,
    fit,
    positioning,
    shapes.geometric
}

\setmainfont{Times New Roman}
\setsansfont{Noto Sans SC}
\setCJKmainfont{Noto Sans SC}
\setCJKsansfont{Noto Sans SC}

\definecolor{ink}{HTML}{1F2937}
\definecolor{muted}{HTML}{64748B}
\definecolor{lineblue}{HTML}{2563EB}
\definecolor{prefill}{HTML}{E7F0FF}
\definecolor{prefillLine}{HTML}{2F5597}
\definecolor{decode}{HTML}{EAF7EA}
\definecolor{decodeLine}{HTML}{2F7D32}
\definecolor{kv}{HTML}{FFF3D8}
\definecolor{kvLine}{HTML}{B7791F}
\definecolor{accent}{HTML}{FDE8E8}
\definecolor{accentLine}{HTML}{C2410C}
\definecolor{panel}{HTML}{F8FAFC}
\definecolor{panelLine}{HTML}{94A3B8}
\definecolor{violet}{HTML}{F1ECFF}
\definecolor{violetLine}{HTML}{6D28D9}
\definecolor{teal}{HTML}{DFF7F3}
\definecolor{tealLine}{HTML}{0F766E}

\tikzset{
    every picture/.style={font=\sffamily},
    title/.style={font=\sffamily\bfseries\Large, text=ink},
    subtitle/.style={font=\sffamily\small, text=muted, align=center},
    block/.style={
        draw=panelLine,
        rounded corners=5pt,
        fill=white,
        line width=0.8pt,
        align=center,
        inner xsep=7pt,
        inner ysep=6pt
    },
    pblock/.style={block, fill=prefill, draw=prefillLine},
    dblock/.style={block, fill=decode, draw=decodeLine},
    kvblock/.style={block, fill=kv, draw=kvLine},
    warnblock/.style={block, fill=accent, draw=accentLine},
    vblock/.style={block, fill=violet, draw=violetLine},
    tblock/.style={block, fill=teal, draw=tealLine},
    cluster/.style={
        draw=panelLine,
        rounded corners=6pt,
        fill=panel,
        line width=0.8pt,
        inner sep=8pt
    },
    arrow/.style={-{Latex[length=3mm,width=2mm]}, draw=ink, line width=0.85pt},
    bluearrow/.style={arrow, draw=prefillLine},
    greenarrow/.style={arrow, draw=decodeLine},
    kvarrow/.style={arrow, draw=kvLine},
    orangearrow/.style={arrow, draw=accentLine},
    control/.style={-{Latex[length=2.6mm,width=1.8mm]}, draw=muted, line width=0.75pt, dashed},
    label/.style={font=\sffamily\scriptsize, text=muted, align=center},
    small/.style={font=\sffamily\scriptsize, align=center},
    rank/.style={
        draw=panelLine,
        rounded corners=3pt,
        fill=white,
        minimum width=12mm,
        minimum height=6mm,
        font=\sffamily\scriptsize,
        align=center
    },
    head/.style={
        draw=ink,
        line width=0.35pt,
        minimum width=6.5mm,
        minimum height=4.8mm,
        font=\sffamily\tiny,
        align=center
    },
    metric/.style={
        draw=panelLine,
        rounded corners=4pt,
        fill=white,
        font=\sffamily\scriptsize,
        align=center,
        inner xsep=5pt,
        inner ysep=4pt
    }
}


\begin{document}
\begin{tikzpicture}[x=1cm,y=1cm]
    \node[title] at (8.6,8.1) {研究路线与实验闭环};
    \node[subtitle] at (8.6,7.66) {按“性能剖析、PD/KV 适配、Decode 动态 CP、运行时弹性与 KV 去冗余”逐步推进，并用端到端指标反向校准设计};

    \node[cluster, minimum width=3.3cm, minimum height=4.7cm] (input) at (1.95,4.45) {};
    \node[font=\sffamily\bfseries\small] at (1.95,6.35) {实验输入};
    \node[metric, fill=white, minimum width=2.55cm] at (1.95,5.55) {模型：Qwen / DeepSeek};
    \node[metric, fill=white, minimum width=2.55cm] at (1.95,4.82) {平台：昇腾 910 系列};
    \node[metric, fill=white, minimum width=2.55cm] at (1.95,4.09) {上下文：4K--1M};
    \node[metric, fill=white, minimum width=2.55cm] at (1.95,3.36) {负载：1--16 rps};
    \node[metric, fill=white, minimum width=2.55cm] at (1.95,2.63) {对照：混合部署/静态 CP};

    \node[pblock, minimum width=2.55cm, minimum height=2.0cm] (s1) at (5.15,5.65) {\textbf{1 性能剖析}\\[2pt]TTFT/TPOT/ITL\\HBM 与 AI Core};
    \node[kvblock, minimum width=2.55cm, minimum height=2.0cm] (s2) at (8.15,5.65) {\textbf{2 PD/KV 适配}\\[2pt]TP/PP 映射\\KV layout 转换};
    \node[dblock, minimum width=2.55cm, minimum height=2.0cm] (s3) at (11.15,5.65) {\textbf{3 动态 CP}\\[2pt]Domain 调度\\跨 DP attention};
    \node[vblock, minimum width=2.55cm, minimum height=2.0cm] (s4) at (14.15,5.65) {\textbf{4 弹性与去冗余}\\[2pt]DP 扩缩容\\TP 域 DCP};

    \node[small, text=muted] at (5.15,4.17) {2026.06--08};
    \node[small, text=muted] at (8.15,4.17) {2026.09--11};
    \node[small, text=muted] at (11.15,4.17) {2026.12--2027.01};
    \node[small, text=muted] at (14.15,4.17) {2027.02--03};

    \draw[bluearrow] (input.east) -- (s1.west);
    \draw[bluearrow] (s1.east) -- (s2.west);
    \draw[kvarrow] (s2.east) -- (s3.west);
    \draw[greenarrow] (s3.east) -- (s4.west);

    \node[cluster, minimum width=11.6cm, minimum height=1.7cm] (eval) at (9.65,2.15) {};
    \node[font=\sffamily\bfseries\small] at (4.65,2.55) {统一评估};
    \node[metric, fill=white, minimum width=3.15cm] at (6.95,2.55) {延迟：TTFT / TPOT / ITL};
    \node[metric, fill=white, minimum width=3.35cm] at (10.35,2.55) {资源：显存峰值 / HBM 读写};
    \node[metric, fill=white, minimum width=3.3cm] at (13.72,2.55) {系统：QPS / 负载均衡};
    \node[metric, fill=white, minimum width=7.45cm] at (10.35,1.72) {补充指标：带宽利用率、KV 传输时间、layout 转换时间、输出一致性};

    \draw[control] (s1.south) -- (eval.north -| s1.south);
    \draw[control] (s2.south) -- (eval.north -| s2.south);
    \draw[control] (s3.south) -- (eval.north -| s3.south);
    \draw[control] (s4.south) -- (eval.north -| s4.south);
    \draw[control] (eval.west) to[out=180,in=-90] node[label, left] {消融与回归} (s1.south);

    \node[warnblock, minimum width=3.2cm, minimum height=4.7cm] (output) at (17.1,4.45) {};
    \node[font=\sffamily\bfseries\small, text=accentLine] at (17.3,6.35) {预期产出};
    \node[small] at (17.3,5.58) {昇腾 PD 分离\\KV 传输与转换方法};
    \node[small] at (17.3,4.78) {Decode 动态 CP\\调度原型与实验};
    \node[small] at (17.3,3.98) {DP 动态扩缩容\\TP 域 KV 去冗余};
    \node[small] at (17.3,3.18) {论文图表、消融分析\\答辩材料};

    \draw[orangearrow] (s4.east) -- (output.west);
    \draw[orangearrow] (eval.east) to[out=0,in=-90] (output.south);
\end{tikzpicture}
\end{document}
